% ============================================================
% APÉNDICE A - SOFTWARE DESARROLLADO
% ============================================================
\chapter{Software desarrollado}
\label{ap:software}

\section{Introducción}
Como parte de este Trabajo Fin de Máster se ha desarrollado \textit{Praeco}, un sistema software orientado a facilitar el acceso a la información textual presente en el etiquetado de productos mediante técnicas de visión artificial, reconocimiento óptico de caracteres y síntesis de voz. La versión final se organiza en tres aplicaciones independientes: \textbf{Praeco Desarrollo}, \textbf{Praeco Usuario} y \textbf{Praeco Análisis}.

\section{Repositorio y versión de referencia}
El código fuente correspondiente a la versión final de Praeco se encuentra disponible públicamente en el repositorio GitHub del proyecto:
\begin{center}
\url{https://github.com/fjlopezlp-beep/Asistente-Discapacidad-Visual}
\end{center}
La versión utilizada como referencia para la entrega final se identifica mediante la etiqueta \texttt{v2.3.1-TFM} y la \textit{release} \textbf{Praeco v2.3.1 -- TFM}. Las tres aplicaciones se identifican con la versión \texttt{2.3.1+1}.

\section{Estructura del repositorio}
El repositorio separa las tres aplicaciones en los directorios \texttt{Praeco\_Desarrollo}, \texttt{Praeco\_Usuario} y \texttt{Praeco\_Analisis}. Cada aplicación conserva su propio código fuente, recursos, configuración Android y ficheros \texttt{pubspec.yaml} y \texttt{pubspec.lock}. No se incluyen cachés de Flutter, resultados de compilación ni configuraciones locales del entorno de desarrollo.

\section{Praeco Desarrollo}
Praeco Desarrollo corresponde a la aplicación utilizada durante las fases de implementación, aprendizaje, experimentación y validación. Integra la adquisición y evaluación de calidad de la imagen, selección de la región de interés, OCR mediante Google ML Kit, filtrado y reconstrucción del texto, corrección léxica, recursos de aprendizaje adaptativo, detección de términos asociados a alérgenos, síntesis de voz y evaluación mediante CER y WER. Durante la validación final, los recursos de aprendizaje se mantienen congelados.

\section{Praeco Usuario}
Praeco Usuario constituye la versión destinada al funcionamiento normal del asistente. No requiere la introducción de \textit{ground truth} ni realiza el cálculo de CER y WER durante su utilización. Emplea los recursos lingüísticos previamente obtenidos y congelados, manteniendo estable el comportamiento del corrector.

\section{Praeco Análisis}
Praeco Análisis es una aplicación independiente destinada al tratamiento estadístico de los resultados experimentales. Permite estudiar CER, WER, comportamiento del corrector, variables de calidad de imagen y tiempos de procesamiento. No necesita los diccionarios empleados por el corrector, puesto que analiza resultados previamente obtenidos y no ejecuta de nuevo el proceso de corrección.

\section{Requisitos de compilación}
Cada aplicación dispone de su propio fichero \texttt{pubspec.yaml}. Una vez instalado y configurado Flutter, las dependencias pueden recuperarse mediante \texttt{flutter pub get}; la aplicación puede ejecutarse con \texttt{flutter run} y puede generarse un APK de distribución mediante \texttt{flutter build apk --release}.

\section{Reproducibilidad}
La combinación del código fuente, los recursos necesarios, la definición y bloqueo de dependencias, la etiqueta Git específica y la \textit{release} asociada permite identificar de forma inequívoca el estado del software utilizado como referencia para el Trabajo Fin de Máster.

\section{Datos generados durante la experimentación}
Durante la ejecución de Praeco Desarrollo se utiliza el directorio \texttt{TFM\_OCR} del dispositivo Android para almacenar datos generados durante los experimentos. Este directorio constituye un espacio de trabajo y no forma parte del código fuente necesario para compilar la aplicación.

Para garantizar la trazabilidad de la evaluación presentada en esta memoria se ha preparado un conjunto de validación independiente, publicado como archivo adicional de la \textit{release} \texttt{v2.3.1-TFM}. El archivo \texttt{Praeco\_Dataset\_Validacion\_TFM.zip} contiene el fichero de resultados correspondiente a 63 registros de validación con el aprendizaje desactivado, las imágenes asociadas y la información de trazabilidad disponible.

\section{Disponibilidad del software}
La versión de referencia utilizada para la entrega de este Trabajo Fin de Máster es \textbf{Praeco v2.3.1 -- TFM}, identificada mediante la etiqueta \texttt{v2.3.1-TFM} en el repositorio indicado anteriormente.
